\section{Исследование и построение решения задачи}
\label{sec:Chapter3} \index{Chapter3}

    Рассмотрим некоторые теоретические аспекты, которые требуются для выполнения
задач данной работы.
\section{Устройство верификатора}

\subsection{Режимы верификации:}
Верификация класса может начинаться в различные моменты времени, в зависимости от целей которые преследуются:

\begin{enumerate}
    \item \textbf{Верификация времени установки}
    Всё приложение, а соответственно его классы, проходят проверку на этапе установки на устройство. Данный 
    подход хорощ, когда мы хотим исключить довольно дорогостоящие проверки из исполнения, ускорить старт приложения,
    возможность агрессивной компиляции в машинный код, всё перечисленное положительно сказывается на исполнении. Однако хоть такой подход и является стандартом индстрии, ультимативным он не является, т.к. имеет ряд 
    недостатков:в данном режиме отсутствует информация какие классы реально будут загружены, какие ветки будут выполнены, 
    какие типы появятся.Но самой главной проблеиой данного подхода является то, что он плохо масштабируется по количеству изменений. 
    Рассмотрим пример: у нас было несколько файлов, они были проверифицированы и откампилированы в OAT файл 
    с машинным кодом, после были внесены какие-то изменения в исходный код, а так OAT файлы являются цельными
    артифактами, то они чувствительны к малейшим изменениям, которые могут сломать метоинформацию, то нам 
    приходится перекомпилировать и проводить верификацию, что является довольно энергоемким процессом. 
    
    \item \textbf{Верификация времени линковки}
    При данном подходе класс будет верифицирован во время его подготовки к исполнению, которая может быть вызвана, 
    например, созданием его экземпляра. Но т.к. во время верификации будут рекурсивно загружены все зависимые классы, что
    может привести к деоптимизациям, связанным с девиртуализацией(например например Оптимизация иерархи класса(cha)).
    \item \textbf{Верификация "на лету"} 
    Данный подход является наиболее гибким, верифиция будет вызвана только при первом вызове метода. Так же данный подход
    является более инкрементальным и изменения проведенные в конкретном методе приведут только к верификации данного метода
    или его класса, такой подход более подходящий для jit.  
\end{enumerate}

\subsection{Cхема работы: }
Сначала рассммотрим устройство верификатора целевой платформы. Единицей верификации
является метод. Для каждого верифицируемого метода создается экземпляр объекта задачи для верифицирования(Job). 
Процесс создания задачи состоит из разрешение типов входных параметров и возвращаемого типа, построения
графа потока управления, а также создния верификаторного контекста, который в себе содержит содержит контекст 
регистров, который является одниим из главных используемых инструментов во время посимвольного исполнения. Разберем
каждый из этапов

\subsubsection{Разрешение идентификаторов}
\subsubsection{Разрешение типов}
\subsubsection{Проверка графа потока управления} % TODO: Mostly AI GENERATED => recheck, refactor, enchance...

Верификатор исследуемой виртуальной машины выполняет большинство проверок графа потока управления, которые были
перечислены ранее: проверка на то, что адрес перехода - начало инструции, проверка отсутствия переходов за пределы
метода и т.д. При этом структура метода имеет свою специфику, соответсвие которой должно проверяться
верификатором: последовательность блоков кода, перемежающихся 
с обработчиками исключений, поддерживая даже вложенные конфигурации обработчиков исключений, когда внутренний 
обработчик находится внутри внешнего. Однако эта гибкость уравновешивается строгим набором правил, 
регулирующих переходы потока управления.

\begin{itemize}
    % \item Выполнение за пределами тела метода: Любая форма прямой передачи потока управления за пределы тела метода, такая как явный переход или неявное проваливание в его конце, запрещена. Это гарантирует, что методы завершаются через предусмотренные механизмы возврата или распространения исключений. Условные переходы в конце метода также в настоящее время не допускаются из-за ограничений в способности статического верификатора окончательно доказать их безопасность.
    % \item Переходы из кода в обработчик исключений: Прямые переходы (jmp) или проваливания из обычных сегментов кода непосредственно в точку входа обработчика исключений не допускаются. Стандартным и предпочтительным механизмом передачи управления обработчику исключений является исключительно инструкция throw. Исключение из этого поведения по умолчанию может быть разрешено через опцию C-TO-H, что указывает на преднамеренное изменение стандартной строгости. Более того, переходы из обычного кода в произвольную точку внутри тела обработчика исключений (а не только в его начало) также строго запрещены, что обеспечивает чистую точку входа для обработки исключений.
    % \item Переходы из обработчика в обработчик: Поток управления непосредственно из одного обработчика исключений в другой — будь то прямой переход или неявное проваливание — запрещен. Это правило предотвращает сложные и потенциально неуправляемые цепочки обработки исключений, способствуя более изолированной и предсказуемой логике восстановления в каждом обработчике. Аналогично, переходы внутрь тела другого обработчика изнутри другого обработчика, или проваливания из внутреннего обработчика в код его внешнего обработчика, также ограничены.
    \item Переходы из обработчика в код: Безусловные переходы из обработчика исключений обратно в обычный сегмент кода, запрещены. Это предотвращает произвольное возобновление нормального выполнения обработчиками исключений, вместо этого направляя их либо на повторное возбуждение исключения, явное завершение метода, либо на тщательно определенный путь восстановления, который не входит в основной поток кода в произвольной точке.
\end{itemize}


\begin{figure}[H] % TODO: Too many different styles of graphs
    %     =========
    %     ...
    %     code
    %     ...
    %     lbl:   <-------\
    %     ...            |
    %     ---------      |
    %     ...            |
    %     exception      |
    %     handler        |
    %     ...            |
    %     jmp lbl  ------/
    %     ...
    %     ---------
    %     ...
    \centering
    \includegraphics[width=0.2\textwidth]{"./images/illegal_jmp.png"} % Путь к файлу
    \caption{Пример некорректного перехода.}
    \label{fig:type_system}
\end{figure}


\subsubsection{Абстрактная интерпритация}
Абстрактная интерпретация представляет собой формальный метод статического анализа программ,
предназначенный для аппроксимации их поведения с целью доказательства корректности определённых свойств без 
выполнения программы на конкретных входных данных.

Метод был предложен Patrick Cousot и Radhia Cousot в 1970. Абстрактная интерпретация позволяет построить корректное 
приближение семантики программы, охватывающее все возможные её исполнения.

В процессе данного анализа происходит работа не с конкретными значениями, а с типами. 
% TODO: Мб это куда-то вынести, но пока хз пусть тут побудет 
При этом верификаторная система 
типов может отличается от системы типов среды исполнения, ниже приведен пример системы типов целевой платформы.
При этом отношение прототипирования для классов, определяемых пользователем, осуществляется по средствам механизмов
среды исполнения. 

\begin{figure}[H] % TODO: Мб 'H' плохо из-за нарушения ГОСТ
    \centering
    \includegraphics[width=0.8\textwidth]{"./images/type_system.png"} % Путь к файлу
    \caption{Типовая система верификатора исследуемой платформы.}
    \label{fig:type_system}
\end{figure}

исследуемой и используя их свойства, строится
некое приближение исполнения реальной программы, при этом данное приближение должно гаранитировать 
корректность реально исполняемой программы. Данный алгоритм должен обладать следующими свойствами:
\textit{достоверность} - если верификатор успешно отработал, то ошибка времени исполнения не может возникнуть,
\textit{неполнота} - т.е. верификатор может отвергнуть валидный код. Разберем второе свойство более подробно. Кроме того 
анализ должен быть конечным

% \subsubsubsection{Конкретная и абстрактная семантика}

Пусть программа $P$ обладает конкретной семантикой, которая описывает точное поведение программы при всех возможных 
входных данных. Конкретная семантика, как правило, оперирует бесконечными множествами состояний, что делает прямой анализ вычислительно неразрешимым, т.к 
количество возможных состояний растет экспоненциально. Для решения данной проблемы вводится абстрактная семантика, которая:
оперирует элементами абстрактного домена, является конечной, корректной аппроксимацией конкретной семантики.
Между конкретной и абстрактной семантиками определяются функции:
\begin{itemize}
    \item функция абстракции $\alpha : S_c \rightarrow S_a$;
    \item функция конкретизации $\gamma : S_a \rightarrow S_c$.
\end{itemize}
образующие галуасово соединение. По определению галуасово соединение это:
$$[\forall c\in C \;и\; \forall a \in A, \text{уде A и C некие абстакный и конкретный домены}] \longmapsto $$
$$ \longmapsto [\alpha(c) \sqsubseteq_A a \Longleftrightarrow c \sqsubseteq_C \gamma(c)] $$
Выполнения данного условия является достаточным гарантирования корректности верификации. Рассмотрим, как данная теория 
отражена на практике. 

Примерами \textit{абстрагирования} могут послужить операции \textit{соединения(join)} и \textit{расширения(widennig)}
\begin{itemize}
    \item Оператор \textit{соединения} ($\sqcup$) представляет собой наименьшую
        верхнюю грань в абстрактном домене. Он
        заключается в объединении информации из различных путей выполнения
        программы для обеспечния наименее грубой аппроксимации,
        покрывающей несколько состояния. В алгоритмах абстрактной интерпретации join
        применяется в точках слияния графа потока управления. Програмно это может выражаться в следующем: в одной ветке
        регистр $v_1$ имеет тип $String$, а по другой тип $int$. В итоге после слияния может произойти следующее:
        $$String \sqcup Int = Object$$, если типовая система верификатора(которая в общем и конкретно нас интересующем 
        случае не равна типовой системе виртуальной машины) подразумевает, что $Object$ — общий супертип(здесь и далее под 
        супертипом подразумевается базовый класс). Или же возможно образования специального типа \textit{объединение(union)}
        $$String \sqcup Int = (String | Int)$$
    \item Для обеспечения завершения абстрактной интерпретации при анализе программ
        с циклами используется оператор \textit{расширения}($\triangledown$). Расширение представляет собой
        специальную бинарную операцию над элементами абстрактного домена,
        ускоряющую достижение неподвижной точки путём ослабления точности
        аппроксимации. Применение расширения гарантирует завершение анализа за
        конечное число шагов, что является критически важным свойством для
        практических верификаторов байт-кода. Рассмотрим пример, когда используется данный подход:
    
\begin{figure}[H]
\centering
\begin{tikzpicture}[node distance=2cm and 3cm]

\node (entry) [start] {Entry};
\node (x0) [process, below of=entry] {x = 0};
\node (H) [decision, below of=x0, yshift=-0.5cm] {x < 100 ?};
\node (xinc) [process, right of=H, yshift=-3cm, xshift=-2cm] {x = x + 1};
\node (stop) [stop, below of=H, yshift=-4cm] {Stop};

% Стрелки
\draw [arrow] (entry) -- (x0);
\draw [arrow] (x0) -- (H);
\draw [arrow] (H.east) -- ++(0.7cm,0) -- ++(0,-5cm) -- ++(-2.3cm,0)-- (stop.north);
\draw [arrow] (H.south) -- (xinc.north);
\draw [arrow] (xinc.south) -- ++(0,-1cm) --  ++(-2cm,0) -- ++(0,4.5cm) -- (H.west);

% Метки стрелок
\node [right=1.5cm] at ($(H.east)!0.5!(stop.north)$) {No};
\node [left=0.1cm] at ($(H.south)!0.5!(xinc.north)$) {Yes};

\end{tikzpicture}
\caption{Блок-схема цикла с условием для демонстрации операции \textit{расширения}.} 
\label{fig:cycle}
\end{figure}
\end{itemize}

Используем интервальный домен:
$$x\in[l,r]$$
$$[l_1, r_1] \sqsubseteq [l_1, r_2] \Leftrightarrow l_2 \leq l_1 \wedge  r_1 \leq r_2 $$
На обратной дуге мы заменили возрастающую цепочку принимаемых $x$ значений:
$$[0,0]\sqsubseteq[0,1]\sqsubseteq[0,2] \sqsubseteq ...$$

Данный цепочка может возрастать долго, поэтому домен состояний $x$ будет определен, как $[0,+\inf]$
\subsection{ets opcodes}
\enlargethispage{\baselineskip}

В исследуемой платформе часть семантики языка ETS не может быть естественно выражена только с помощью
базового набора объектных инструкций наподобие обычных загрузок полей, вызовов виртуальных методов и
сравнений ссылок. По этой причине в ISA выделено отдельное ETS-расширение с собственным префиксом
инструкций. Такое расширение позволяет сделать байт-код ближе к исходной семантике языка и явно
зафиксировать в коде те операции, которые зависят не только от низкоуровневого представления объектов,
но и от языковых правил ETS.

\subsubsection{Мотивация появления ETS-инструкций}

Одной из отличительных особенностей целевой платформы по сравнению с классическими java-подобными
виртуальными машинами является поддержка объединения типов. Под объединением типов в данной работе
понимается ссылочный тип, принимающий значения одного из нескольких классов. Для такого значения
необходимо уметь безопасно обращаться к общим членам объединения.

\begin{lstlisting}[language=Python, label={lst:union_example}, caption={Пример кода с объединением типов}]
    class A { foo() { return 1 }}
    class B { foo() { return 2 }}
    
    function bar(v : A|B) { 
        v.foo()
    }
\end{lstlisting}

\begin{lstlisting}[language=Python, label={lst:union_bc}, caption={Представление кода из Листинга~\ref{lst:union_example} на уровне байткода}]
    lda.obj a0
    sta.obj v1
    ets.call.name.short union.%%union_prop-A|B.foo:(union.%%union_prop-A|B), v1
    return.void
\end{lstlisting}

Компилятор генерирует синтетическое описание объединения и на этапе кодогенерации строит инструкцию
обращения по имени к общему члену объединения. Важно, что язык не допускает создания объекта
объединения через оператор \texttt{new}. Следовательно, во время исполнения в регистре находится
экземпляр конкретного класса, например \texttt{A} или \texttt{B}, а инструкция задает не адрес
конкретного слота, а имя требуемого поля или метода. Таким образом, ETS-инструкции позволяют
перенести операцию разрешения члена на фактический тип объекта во время исполнения.

\subsubsection{Основные семейства ETS-инструкций}

Содержательно ETS-расширение можно разбить на несколько семейств:
\begin{itemize}
    \item инструкции чтения поля по имени:
    \texttt{ets.ldobj.name}, \texttt{ets.ldobj.name.64}, \texttt{ets.ldobj.name.obj};
    \item инструкции записи поля по имени:
    \texttt{ets.stobj.name}, \texttt{ets.stobj.name.64}, \texttt{ets.stobj.name.obj};
    \item инструкции вызова метода по имени: \texttt{ets.call.name.short}, \texttt{ets.call.name},
    \texttt{ets.call.name.range};
    \item специальные инструкции для работы с языковым значением \texttt{nullvalue}:
    \texttt{ets.ldnullvalue}, \texttt{ets.movnullvalue}, \texttt{ets.isnullvalue};
    \item инструкции сравнения и определения типа:
    \texttt{ets.equals}, \texttt{ets.strictequals}, \texttt{ets.typeof}.
\end{itemize}

Первое семейство отвечает за чтение и запись поля объекта по имени. По сравнению с обычными
\texttt{ldobj}/\texttt{stobj} здесь важен не только идентификатор в constant pool, но и то, что
итоговое разрешение выполняется относительно фактического класса объекта-получателя. Существование
вариантов \texttt{.64} и \texttt{.obj} отражает различие в представлении результата: отдельно
обрабатываются 64-битные примитивы и ссылочные значения, тогда как базовый вариант используется
для 32-битных значений.

С точки зрения реализации такая инструкция интересна тем, что рантайм сначала пытается найти реальное
поле с требуемым именем в классе объекта и его базовых классах. Если подходящее поле отсутствует,
то доступ может быть интерпретирован как вызов аксессора. Иными словами, операция по имени в ETS
абстрагируется от конкретной формы представления члена: логически это может быть либо поле, либо
getter/setter-метод с согласованным именем. Такой подход позволяет сохранить естественную объектную
модель языка на уровне байткода, не сводя все обращения исключительно к полям.

Второе семейство, \texttt{ets.call.name.*}, предназначено для вызова метода по имени. В отличие
от обычного виртуального вызова, где компилятор чаще всего уже знает точную сигнатуру члена,
здесь существенен сам факт вызова по имени относительно текущего объекта. Формы \texttt{short},
\texttt{full} и \texttt{range} различаются способом передачи аргументов: короткая форма удобна
для небольшого числа регистров, а range-вариант используется, когда аргументов больше и они
расположены в непрерывном диапазоне регистров.

Отдельного упоминания заслуживает то, что ETS-инструкции доступа по имени не являются исключительно
механизмом для объединений типов. Они задают обобщенный способ работы с объектной семантикой ETS,
когда конкретный член определяется именем и текущим фактическим классом объекта. Это делает их
полезными и в иных языковых конструкциях, где требуется отложенное разрешение члена.

\subsubsection{Специальные семантики ETS}

Помимо доступа к полям и вызовов, ETS-расширение содержит инструкции, отражающие специальные
языковые понятия. Семейство \texttt{nullvalue} включает три инструкции:
\texttt{ets.ldnullvalue}, \texttt{ets.movnullvalue}, \texttt{ets.isnullvalue}. Они работают со
специальным ссылочным значением \texttt{nullvalue}, которое используется для представления
nullish-состояния в семантике языка. Наличие отдельного семейства инструкций показывает, что для
платформы важно различать обычные объектные ссылки и специальное пустое значение, значимое на
уровне языка.

\begin{samepage}
Отдельную группу составляют инструкции сравнения. \texttt{ets.equals} и
\texttt{ets.strictequals} реализуют две разновидности сравнения. Такое разделение делает байт-код
более выразительным: различие между «мягким» и «строгим» сравнением не скрывается внутри
библиотечного вызова, а явно фиксируется на уровне ISA. Инструкция \texttt{ets.typeof}, в свою
очередь, инкапсулирует получение строкового описания типа значения в терминах ETS.

\par\noindent
В итоге ETS-расширение можно рассматривать как промежуточный слой между общей объектной моделью
виртуальной машины и высокоуровневой семантикой языка. Оно позволяет сохранить в байт-коде сведения
о том, что операция возникла именно из ETS-конструкции, а не была сведена к набору универсальных
низкоуровневых действий.
\end{samepage}

\subsection{any opcodes}

Другой важной особенностью исследуемой платформы является поддержка взаимодействия со средой
динамически типизированных языков. Для ArkTS это особенно актуально в сценариях межъязыкового
взаимодействия, когда статически типизированный код вызывает функции динамического мира, получает
из него значения и затем обращается к их свойствам или методам.

\subsubsection{Модель значения \texttt{any}}

В ISA платформы для таких сценариев вводится специальное значение \texttt{any}. Существенно, что
\texttt{any} не является просто синонимом ссылочного типа \texttt{Object}. Динамическое значение
хранит не только само содержимое, но и информацию о его runtime-типе. В \texttt{any} может быть
упаковано как примитивное значение, так и объектная ссылка. Благодаря этому виртуальная машина
может единообразно представлять результаты вычислений в динамической среде.

Важным следствием такой модели является то, что регистры и аккумулятор должны быть способны хранить
значения \texttt{any}. Поэтому в наборе инструкций присутствуют не только операции доступа к полям
и вызова методов, но и более базовые операции перемещения динамических значений между регистрами,
аккумулятором и непосредственными литералами.

\subsubsection{Базовые динамические инструкции}

Поддержка \texttt{any} начинается с семейства \texttt{dyn}-инструкций. К ним относятся
\texttt{mov.dyn}, \texttt{lda.dyn}, \texttt{sta.dyn}. Отдельно выделяются инструкции загрузки
литералов \texttt{ldai.dyn} и \texttt{fldai.dyn}. Их назначение состоит в том, чтобы перемещать
динамические значения между регистрами и аккумулятором и загружать в вычисление непосредственные
константы, уже представленные в форме \texttt{any}.

Архитектурно это означает, что граница между статической и динамической частями программы видна уже
на уровне байткода. Компилятор не вынужден неявно преобразовывать все значения к универсальному
объектному типу. Вместо этого факт перехода к динамической модели фиксируется отдельным набором
инструкций.

\subsubsection{Доступ к свойствам и элементам динамических значений}

Непосредственное взаимодействие с динамическими объектами осуществляется с помощью префиксного
семейства \texttt{any.*}.

\begin{lstlisting}[language=Python, label={lst:any_example}, caption={Пример взаимодействия со значением из динамической среды}]
    // dynamic.js
    class A { prop = 1}
    function foo() { return new A() }

    // static.ets
    import { A, foo } from "./dynamic.js"
    let a = foo()["prop"]
\end{lstlisting}

\begin{lstlisting}[language=Python, label={lst:any_bc}, caption={Представление кода из Листинга~\ref{lst:any_example} на уровне байткода}]
    lda.obj a0
    sta.obj v1
    any.ldbyname v1, "prop"
    return.obj
\end{lstlisting}

Операции доступа можно разделить по способу адресации:
\begin{itemize}
    \item чтение по имени свойства: \texttt{any.ldbyname}, \texttt{any.ldbyname.v};
    \item запись по имени свойства: \texttt{any.stbyname}, \texttt{any.stbyname.v};
    \item по числовому индексу: \texttt{any.ldbyidx}, \texttt{any.stbyidx};
    \item по произвольному значению-ключу: \texttt{any.ldbyval}, \texttt{any.stbyval}.
\end{itemize}

Такое деление хорошо соответствует операциям динамических языков. Доступ по имени моделирует
обращение к свойству объекта через строковый идентификатор. Доступ по индексу естественным образом
подходит для массивоподобных структур. Наконец, доступ по значению-ключу отражает общую форму
индексатора, где в качестве ключа может выступать произвольное значение.

Стоит заметить, что часть инструкций записывает результат в аккумулятор, а часть --- непосредственно
в регистр. Наличие двух вариантов делает ISA более гибким и позволяет избегать лишних промежуточных
перемещений значения, если следующий этап вычисления ожидает аргумент именно в регистре.

\subsubsection{Динамические вызовы}

Помимо доступа к полям и индексаторам, семейство \texttt{any.*} покрывает основные формы вызова,
характерные для динамического кода:
\begin{itemize}
    \item вызов функции: \texttt{any.call.0}, \texttt{any.call.short}, \texttt{any.call.range};
    \item вызов метода объекта по имени:
    \texttt{any.call.this.0}, \texttt{any.call.this.short}, \texttt{any.call.this.range};
    \item вызов конструктора:
    \texttt{any.call.new.0}, \texttt{any.call.new.short}, \texttt{any.call.new.range};
    \item проверка принадлежности типу: \texttt{any.isinstance}.
\end{itemize}

Здесь снова прослеживается общая архитектурная идея: виртуальная машина не скрывает динамическую
семантику за одним универсальным внешним вызовом, а делает различные сценарии явными. Вызов функции,
вызов метода объекта и конструирование нового объекта выражаются разными байткодами, хотя все они
работают в динамическом контексте. Это облегчает как интерпретацию, так и последующий анализ байткода,
поскольку вид операции известен непосредственно из кода инструкции.

Разделение на формы \texttt{0}, \texttt{short} и \texttt{range} служит той же цели, что и в других
частях ISA: компактное кодирование часто встречающихся случаев при сохранении возможности передавать
произвольное число аргументов. Короткие формы экономят место в байткоде, а range-варианты позволяют
эффективно работать с длинными списками аргументов.

\subsubsection{Роль any-инструкций в архитектуре платформы}

С теоретической точки зрения семейство \texttt{any} важно тем, что оно локализует границу между
статическим и динамическим миром. Вместо попытки представить динамические значения как обычные
ссылки и затем восстанавливать их смысл неявным образом, платформа явно выделяет особый супертип
и специальный набор операций над ним.

Такое решение делает байт-код более выразительным. По последовательности инструкций можно увидеть,
в каких точках программа переходит к динамической модели вычислений, где происходит доступ к свойствам
по имени, где используется индексатор, а где выполняется динамический вызов. Следовательно,
\texttt{any}-байткоды образуют самостоятельный архитектурный слой поддержки межъязыкового взаимодействия
и динамической семантики внутри в остальном статически типизированной платформы.